\documentclass[]{article}
\usepackage{graphicx}
\usepackage[svgnames]{xcolor} 
\usepackage{fancyhdr}

\usepackage[hidelinks]{hyperref}
\usepackage{enumitem}
\usepackage[many]{tcolorbox}
\usepackage{listings }
\usepackage[a4paper, total={6in, 8in}]{geometry}
\usepackage{afterpage}
\usepackage{amssymb}
\usepackage{pdflscape}
\usepackage{textcomp}
\usepackage{xecolor}
\usepackage{rotating}
\usepackage[Kashida]{xepersian}
\usepackage[T1]{fontenc}
\usepackage{tikz}
\usepackage[utf8]{inputenc}
\usepackage{PTSerif} 
\usepackage{seqsplit}
\usepackage{changepage}



\usepackage{listings}
\usepackage{xcolor}
\usepackage{sectsty}
 
\definecolor{codegreen}{rgb}{0,0.6,0}
\definecolor{codegray}{rgb}{0.5,0.5,0.5}
\definecolor{codepurple}{rgb}{0.58,0,0.82}
\definecolor{backcolour}{rgb}{0.95,0.95,0.92}
 
\NewDocumentCommand{\codeword}{v}{
\texttt{\textcolor{blue}{#1}}
}
\lstset{language=java,keywordstyle={\bfseries \color{blue}}}

\lstdefinestyle{mystyle}{
    backgroundcolor=\color{backcolour},   
    commentstyle=\color{codegreen},
    keywordstyle=\color{magenta},
    numberstyle=\tiny\color{codegray},
    stringstyle=\color{codepurple},
    basicstyle=\ttfamily\normalsize,
    breakatwhitespace=false,         
    breaklines=true,                 
    captionpos=b,                    
    keepspaces=true,                 
    numbers=left,                    
    numbersep=5pt,                  
    showspaces=false,                
    showstringspaces=false,
    showtabs=false,                  
    tabsize=2
}

\lstset{style=mystyle}

 \settextfont[BoldFont={XB Zar bold.ttf}]{XB Zar.ttf}


\setlatintextfont[Scale=0.8,
 BoldFont={LiberationSerif-Bold.ttf}, 
 ItalicFont={LiberationSerif-Italic.ttf}]{LiberationSerif-Regular.ttf}





\newcommand{\inputsample}[1]{
    ~\\
    \textbf{ورودی نمونه}
    ~\\
    \begin{tcolorbox}[breakable,boxrule=0pt]
        \begin{latin}
            \large{
                #1
            }
        \end{latin}
    \end{tcolorbox}
}

\newcommand{\outputsample}[1]{
    ~\\
    \textbf{خروجی نمونه}

    \begin{tcolorbox}[breakable,boxrule=0pt]
        \begin{latin}
            \large{
                #1
            }
        \end{latin}
    \end{tcolorbox}
}

\newenvironment{changemargin}[2]{%
\begin{list}{}{%
\setlength{\topsep}{0pt}%
\setlength{\leftmargin}{#1}%
\setlength{\rightmargin}{#2}%
\setlength{\listparindent}{\parindent}%
\setlength{\itemindent}{\parindent}%
\setlength{\parsep}{\parskip}%
}%
\item[]}{\end{list}}


\definecolor{foldercolor}{RGB}{124,166,198}
\definecolor{sectionColor}{HTML}{ff5e0e}
\definecolor{subsectionColor}{HTML}{008575}


\defpersianfont\titr[Scale=1.5]{Lalezar-Regular.ttf}

\sectionfont{\color{sectionColor}}  % sets colour of sections



\subsectionfont{\color{subsectionColor}}  % sets colour of sections

\renewcommand{\labelitemii}{$\star$}
\renewcommand{\labelitemiii}{$\star\#$}

\begin{document}


%%% title pages
\begin{titlepage}
\begin{center}
        
\vspace{0.7cm}

\includegraphics[width=0.4\textwidth]{sharif1.png}\\
\vspace{0.3cm}
\textbf{ \Huge{\emph درس برنامه‌سازی پیشرفته} }\\
\vspace{0.5cm}
\textbf{ \Large{ فاز صفر پروژه} }
\vspace{0.2cm}
       
 
      \large \textbf{دانشکده مهندسی کامپیوتر}\\\vspace{0.1cm}
    \large   دانشگاه صنعتی شریف\\\vspace{0.2cm}
       \large   ﻧﯿﻢ سال دوم 99-98 \\\vspace{0.10cm}
      \noindent\rule[1ex]{\linewidth}{1pt}
اساتید:\\
    \textbf{{مهدی مصطفی‌زاده، ایمان عیسی‌زاده، امیر ملک‌زاده، علی چکاه}}



    \vspace{0.20cm}

   مهلت ارسال:\\
    \textbf{{12 فروردین}}\\
    \textbf{{ساعت 23:59:59}}

    \vspace{0.10cm}
مسئول پروژه:\\
    \textbf{{احمد سلیمی}}
    
        \vspace{0.10cm}
مسئول فاز اول:\\
    \textbf{{تست}}
    
        \vspace{0.10cm}
طراحان فاز اول:\\
    \textbf{{تست}}
    
        \vspace{0.10cm}
مسئول تیم تنظیم داک:\\
    \textbf{{تست}}
    
            \vspace{0.10cm}
اعضای تیم تنظیم داک:\\
    \textbf{{تست}}
\end{center}
\end{titlepage}
%%% title pages


%%% header of pages
\newpage
\pagestyle{fancy}
\fancyhf{}
\fancyfoot{}
\cfoot{\thepage}
\lhead{فاز صفر}
\rhead{\includegraphics[width=0.1\textwidth]{sharif.png}\\
دانشکده مهندسی کامپیوتر
}
\chead{پروژه برنامه‌سازی پیشرفته}
%%% header of pages
\renewcommand{\headrulewidth}{2pt}

\KashidaOff


 \Large \textbf{\\\\
به موارد زیر توجه کنید:}

\begin{itemize} 
\item اول \dots{}
     \begin{itemize} 
        \item labelitemii را در فایل پیدا کنید.
        
        \item تست
        \begin{itemize}
        
        \item عبارت labelitemiii را در فایل برای تغییر علامت پشت لیست پیدا کنید.
        
        \item تست کامل
        \end{itemize}
        
     \end{itemize}
\item دوم \dots{}
\item سوم \dots{}
\end{itemize}




\newpage
\section* {{\titr مقدمه - فاز صفر یعنی چی؟}}
در این فاز شما باید مقدمات پروژه را حاضر کنید. این مقدمات شامل ابزارهای مورد استفاده در پروژه و همچنین طراحی معماری پروژه توسط دیاگرام uml است. مراحل این فاز عبارت اند از:

\begin{enumerate}
\item
\hyperref[subsec:trello]{\textcolor{blue}{راه‌اندازی Trello}}

\item
\hyperref[subsec:github]{\textcolor{blue}{راه اندازی مخزن GitHub}}

\item
\hyperref[subsec:uml]{\textcolor{blue}{طراحی UML برای منطق پروژه}}


\end{enumerate}

در بخش بعد، هر یک از این موارد شرح داده شده اند.

\newpage
\section*{{\titr کارهایی که باید در فاز صفر انجام دهید}}


\subsection*{{\titr راه‌اندازی Trello}}
\label{subsec:trello}

مستند آموزشی این بخش، در پیوست این فاز موجود است.

در این پروژه، شما بایستی کارهای هر شخص را در  trello مشخص کرده و پس از انجام آن کار توسط هر شخص، تسک مربوط به آن را  به بخش تسک‌های انجام شده منتقل کنید.

هر یک از اعضای تیم پروژه، باید در سایت trello.com یک حساب کاربری بسازد. سپس یکی از اعضای تیم، یک board جدید ساخته و مابقی اعضا را به آن دعوت کند.

حال در board ساخته شده، تسک‌های زیر را معین کرده، و هر کدام را به یک نفر اختصاص دهید:

\begin{itemize}

\item
ایجاد مخزن github و commit نام و شماره دانشجویی در readme (تسک به نفر اول assign شود)

\item
کامیت نام و شماره دانشجویی در readme (تسک به نفر دوم assign شود)

\item
کامیت نام و شماره دانشجویی در readme (تسک به نفر سوم assign شود)

\end{itemize}

پس از انجام تسک ها آن‌ها را به بخش done منتقل کنید.

در ادامه‌ی کار نیز باید از این ابزار استفاده کنید ‌و تسک‌های خود را مدیریت کنید و در تحویل حضوری هر فاز، بورد شما بررسی می‌شود.


\newpage
\subsection*{{\titr راه‌اندازی مخزن GitHub}}
\label{subsec:github}

همانطور که می‌دانید برای پروژه لازم است با گروهتان بر روی یک مخزن (repository) گیت فعالیت کنید. برای ساختن این مخزن، کافیست وارد
 \href{https://classroom.github.com/g/5Uo3x-M4}{\textcolor{blue}{\underline{این لینک}}} 
 شوید.

ابتدا با لیستی مواجه می‌شوید که شماره دانشجویی تمام افراد در آن موجود است. شماره دانشجویی خود را بیابید و بر روی آن کلیک کنید.

در صفحه‌ی بعد شما باید تیم خود را انتخاب کنید. چنانچه نفر اول گروه خود (سازنده‌ی مخزن) هستید، باید یک تیم بسازید. تنها شماره‌ی گروه پروژه خود را در قسمت نام تیم وارد کنید و تیم را بسازید. نفرات بعدی گروه شما، باید تیم‌شان را از لیست تیم‌های موجود انتخاب کنند و نیازی به ایجاد تیم ندارند.

پس از این مراحل مخزن شما آماده خواهد شد و لینک آن در اختیارتان قرار خواهد گرفت. پس از آماده شدن این مخزن، هر یک از اعضای پروژه باید نام و شماره دانشجویی خود را به فایل readme.md اضافه کند.

\newpage
\subsection*{{\titr طراحی UML برای منطق پروژه}}
\label{subsec:uml}

همزمان با این مستند، مستند فاز اول پروژه نیز منتشر شده است.

برای این بخش، شما لازم است مستند فاز اول پروژه را به طور کامل و دقیق مطالعه نمایید. سپس یک uml مناسب برای تمام بخش‌های منطق پروژه طراحی کنید. فایل ارسالی شما باید یک فایل zip، حاوی تصاویر مربوط به uml باشد.

مستند آموزشی uml نیز در پیوست این فاز موجود است.



\end{document}









